\documentclass[11pt]{article}
\usepackage[margin=1in]{geometry}
\usepackage{amsmath,amssymb,amsthm,mathtools,booktabs}
\usepackage[T1]{fontenc}
\usepackage{lmodern}
\usepackage[hidelinks]{hyperref}
\usepackage{microtype}
\hypersetup{
  pdftitle={Improved lower bound for the Gyarmati-Hennecart-Ruzsa sum-difference constant using masked digits and controlled carries},
  pdfauthor={Logan Kleinwaks}
}

\newtheorem{proposition}{Proposition}
\newtheorem{theorem}[proposition]{Theorem}
\newtheorem{lemma}[proposition]{Lemma}
\newtheorem{remark}{Remark}
\newcommand{\N}{\mathbb{N}_0}
\newcommand{\card}[1]{\lvert #1\rvert}
\newcommand{\angles}[1]{\langle #1\rangle}
\newcommand{\packageurl}{\texttt{https://github.com/kleinwaks/masked-digit-sum-difference-bound}}

\title{Improved lower bound for the\\Gyarmati--Hennecart--Ruzsa
sum-difference constant\\using masked digits and controlled carries}
\author{Logan Kleinwaks}
\date{First version: July 24, 2026; Version 3.2: August 14, 2026}

\begin{document}
\maketitle

\begin{abstract}
Zheng's bounded-digit construction lower bounding the Gyarmati--Hennecart--Ruzsa
sum-difference constant forms integers in a fixed base, restricts every digit
to an interval from zero to a largest digit $B$, and imposes a bound on the
sum of the digits.  We first replace the complete digit interval by an
arbitrary finite mask $M\subseteq\{0,\ldots,B\}$.  When the
base is $2B+1$, addition and subtraction are carry-free, and a
method-of-types argument expresses the resulting limiting bound through two
finite partition polynomials.

We next replace the carry-free base by an arbitrary integer $q>B$.  On the difference
side, each residue block is assigned its exact minimum witness cost; on the
sum side, all representations of the same output block are combined before
they are counted.  Supermultiplicativity and submultiplicativity show that
the corresponding block pressures have limits.  A positive-vector
certificate on a $351951$-state accessible subgraph, together with a
directed $8192$-digit sum-pressure computation, proves
\[
 \boxed{C_{3a}>1.19519192}.
\]
The construction uses $B=26972$, $q=27022$, and
$M=\langle1971,2016,2100,2628,2688,2800\rangle\cap[0,B]$.
Finally, we give a single explicitly specified finite set certifying
$C_{3a}>1.19102809$.  This finite certificate uses neither a
method-of-types limit nor an infinite-block pressure.
\end{abstract}

\section*{Formal verification and AI-use disclosure}
An accompanying Lean~4/\href{https://github.com/leanprover-community/mathlib4}{Mathlib}
formalization was produced with Aristotle (Harmonic) \cite{Aristotle2025} and is included
in the repository.  It machine-checks the definition of $C_{3a}$, the
Gyarmati--Hennecart--Ruzsa finite-set principle used here, the masked-digit
construction, the full controlled finite-block bound, the controlled
difference and sum pressure limits, the finite-state comparison arguments,
the semigroup lemmas, and, in particular, the two conclusions $C_{3a}>1.19519192$ and, from the explicit finite construction, $C_{3a}>1.19102809$.

The source contains no \texttt{sorry}.  Its large exact computations are
evaluated with \texttt{native\_decide}; consequently, the corresponding
axiom report includes \texttt{Lean.trustCompiler} in addition to the standard
Mathlib quotient, choice, and extensionality axioms.

The formal proof of Proposition~\ref{prop:controlled-block} retains the
possible nonzero self-inverse residue and implements the positional
orientation described below by pairing consecutive occurrences from left to
right.  Its block-selection argument uses product-weight concentration with
an even-parity restriction, rather than the fixed exact type displayed in
the paper; it proves the same inequality.  The formal numerical
proof verifies the two pressure inequalities through exact truncated cost
automata with $64003$ difference states and $120001$ sum states.  These are
sound finite relaxations of the exact pressures, but they are not a
line-by-line replay of the independent $351951$-state and $8192$-digit
computations described in Section~\ref{sec:controlled-certificate}.

During this work, the author used several large language models (LLMs).
OpenAI's GPT-5.6 Sol model was used extensively in software development,
analysis of search results, drafting and revising this paper, and preparation
of the verification package.  Google's Gemini 3.1 Pro Preview was used in
software development and analysis of search results; Google's Gemini 3.5
Flash and Gemini 3.6 Flash were used to review the paper; and Anthropic's
Sonnet 5 was used in analysis of search results.  Aristotle produced the
Lean formalization, which was subsequently audited against the mathematical
statements in this paper.  Substantial portions of the initial prose and
source code were generated or revised with this assistance.  The author set
the research goals, directed the computational work, selected which
suggestions to pursue, ran the search and verification software, reviewed the
mathematical arguments and outputs, and takes responsibility for all claims
and references.

\section{Introduction and the finite-set principle}
Let $C_{3a}$ be the supremum of the exponents $\theta$ with the following
property: for every $K>1$, there is a constant $c(K)>0$ and there are pairs
of nonempty finite sets $X,Y\subset\mathbb Z$, for arbitrarily large values of
$\card X$, such that
\[
 \card{X+Y}\le K\card X,
 \qquad
 \card{X-Y}\ge c(K)\card{X+Y}^{\theta}.
\]
Gyarmati, Hennecart, and Ruzsa proved the following
finite-set lower-bound principle \cite[Lemma in Section~2]{GHR2007}: if $U\subset\N$ is finite
and contains zero, and if
\[
 \card{U-U}<2\max U+1,
\]
then
\begin{equation}\label{eq:GHR}
 C_{3a}\ge
 1+\frac{\log\bigl(\card{U-U}/\card{U+U}\bigr)}
          {\log(2\max U+1)}.
\end{equation}

For any finite $U\subset\N$ containing zero, the elementary containment
$U-U\subseteq[-\max U,\max U]$ gives only
$\card{U-U}\le2\max U+1$, whereas the hypothesis of \eqref{eq:GHR} is
strict.  Dilation supplies strictness without requiring any further
information about $U$.  If $\max U>0$ and $V=2U$, then
\[
 \card{V-V}=\card{U-U},\qquad
 \card{V+V}=\card{U+U},\qquad
 \max V=2\max U.
\]
Consequently,
\[
 \card{V-V}=\card{U-U}\le 2\max U+1
 <4\max U+1=2\max V+1.
\]
Thus $V$ satisfies the strict hypothesis of \eqref{eq:GHR}, whether or not
$U$ does, and
\begin{equation}\label{eq:GHR-dilated}
 C_{3a}\ge
 1+\frac{\log\bigl(\card{U-U}/\card{U+U}\bigr)}
          {\log(4\max U+1)}.
\end{equation}
The price is the replacement of $\log(2\max U+1)$ by
$\log(4\max U+1)$.  In the limiting constructions this does not change the
exponential rate; in the finite certificate it adds less than $\log2$ to a
denominator of order $10^9$.  We use the dilation because strictness for the
undilated sets is not otherwise established.

Recent computational work on this constant includes AlphaEvolve
\cite{GeorgievEtAl2025}, Gerbicz's large finite construction
\cite{Gerbicz2025}, and Zheng's large-deviation construction
\cite{Zheng2025}. The \emph{Optimization Constants in Mathematics} repository
\cite{OptimizationConstants2026} subsequently listed the finite-certificate
bounds $1.1740744$ and $1.1835129324$ \cite{Griego2026,Mosaic2026}.
Two submissions merged in July 2026 then gave the limiting values
$1.187326127925948$ \cite{Numaro2026} and $1.19102809$
\cite{Kleinwaks2026}.  The controlled-carry generalization proved here raises
the limiting bound to $1.19519192$, while the separate finite certificate
gives $1.19102809$ without asymptotic analysis.

Gerbicz introduced bounded base digits into the original construction, and
Zheng analyzed the limiting exponential growth of a bounded-digit family.
In Zheng's version the allowed digit alphabet is the complete interval
$\{0,\ldots,B\}$. We replace that interval by a mask $M$. This can remove
many possible sum digits while preserving difference digits $d$ that admit
witnesses $a,b\in M$ with $a-b=d$ and small $a+b$, so that they use little
of the two digit-sum budgets in the construction. The multinomial counting
estimate used below is standard; Proposition~\ref{prop:masked} is the
construction-specific arbitrary-mask consequence.  We first retain the
carry-free base $2B+1$ and then replace it by any $q>B$.  The latter step
allows controlled carries and leads to a transfer operator whose pressure is
the $k\to\infty$ block limit.  After proving and certifying the two bounds, we
give a separate finite-depth witness and then discuss the numerical-semigroup
structure that guided the search.

\begin{theorem}[Main theorem]\label{thm:main}
The Gyarmati--Hennecart--Ruzsa constant satisfies
\[
 C_{3a}>1.19519192.
\]
In addition, a single explicitly specified finite set, without recourse to
either asymptotic theorem in this paper, certifies $C_{3a}>1.19102809$.
\end{theorem}
The two assertions are proved in Sections~\ref{sec:controlled-certificate}
and~\ref{sec:finite-certificate}, respectively.

\section{The masked-digit bound}
Let $M\subset\N$ be finite with $0\in M$ and
$B=\max M\ge1$. Let
\[
 Q=2B+1,
 \qquad S=M+M,
 \qquad D=M-M.
\]
For each $d\in D$, define its minimum witness cost
\begin{equation}\label{eq:kappa}
 \kappa(d)=\min\{a+b:a,b\in M,\ a-b=d\}.
\end{equation}
The function $\kappa$ is even: $\kappa(-d)=\kappa(d)$. Define
\begin{equation}\label{eq:partition-polys}
 P_+(x)=\sum_{s\in S}x^s,
 \qquad
 P_-(x)=\sum_{d\in D}x^{\kappa(d)}.
\end{equation}
We use the following standard lower bound for a multinomial coefficient;
see, for example, Csisz\'ar~\cite[Lemmas~II.1 and~II.2]{Csiszar1998}.
If $E$ is a nonempty finite set and nonnegative integers $m_e$ satisfy
$\sum_{e\in E}m_e=m$, let $\pi_e=m_e/m$. Then
\begin{equation}\label{eq:multinomial-estimate}
 \log\frac{m!}{\prod_{e\in E}m_e!}
 \ge -m\sum_{e\in E}\pi_e\log \pi_e
      -\log\binom{m+\card E-1}{\card E-1},
\end{equation}
where $\log$ is the natural logarithm and $0\log0$ is interpreted as zero.
For fixed $E$, the binomial coefficient in \eqref{eq:multinomial-estimate}
is a polynomial in $m$ of degree $\card E-1$; equivalently,
\[
 \log\binom{m+\card E-1}{\card E-1}=O_E(\log m)=o(m).
\]

\begin{proposition}[Masked-digit bound]\label{prop:masked}
For every $x\in(0,1)$,
\begin{equation}\label{eq:masked-bound}
 C_{3a}\ge
 1+\frac{\log P_-(x)-\log P_+(x)}{\log(2B+1)}.
\end{equation}
Consequently, one may take the supremum of the right-hand side over
$x\in(0,1)$.
\end{proposition}

\begin{proof}
Fix $x\in(0,1)$ and define
\begin{equation}\label{eq:gibbs}
 p_d=\frac{x^{\kappa(d)}}{P_-(x)}
 \qquad(d\in D).
\end{equation}
Since $\kappa$ is even, these weights satisfy $p_d=p_{-d}$. Let
\[
 \mu=\sum_{d\in D}p_d\kappa(d).
\]
The entropy of $p$ is
\[
 G(p)=-\sum_{d\in D}p_d\log p_d
     =\log P_-(x)-\mu\log x.
\]

For each $d\in D$, choose $\alpha_d,\beta_d\in M$ such that
\[
 \alpha_d-\beta_d=d,
 \qquad
 \alpha_d+\beta_d=\kappa(d),
\]
with $(\alpha_{-d},\beta_{-d})=(\beta_d,\alpha_d)$; for $d=0$, take
$(\alpha_0,\beta_0)=(0,0)$.

Let
\[
 R=\log\frac{P_-(x)}{P_+(x)}.
\]
If $R\le0$, then \eqref{eq:masked-bound} follows from the elementary bound
$C_{3a}\ge1$.  For example, fix $0<\delta<K-1$ and take
$X_n=\{0,\ldots,n-1\}$ and
$Y_n=\{0,\ldots,\lfloor\delta n\rfloor-1\}$.  For all sufficiently large
$n$, these intervals satisfy $\card{X_n+Y_n}\le K\card{X_n}$, and their
sumset and difference set have the same cardinality.  Assume henceforth that $R>0$,
and choose
\[
 0<\varepsilon<\frac{R}{-\log x}.
\]
For each positive integer $m$ and each positive $d\in D$, define
\[
 n_d^{(m)}=n_{-d}^{(m)}=\lfloor mp_d\rfloor,
 \qquad
 n_0^{(m)}=m-2\sum_{\substack{d\in D\\d>0}}\lfloor mp_d\rfloor.
\]
These are nonnegative integers for every $m$: indeed, symmetry gives
$2\sum_{d>0}p_d=1-p_0$, and hence
\[
 n_0^{(m)}\ge m-2m\sum_{d>0}p_d=mp_0>0.
\]
Their sum is $m$, and, as $m\to\infty$,
$n_d^{(m)}/m\to p_d$ for every $d\in D$.  Write
$q_d^{(m)}=n_d^{(m)}/m$.  The number of sequences containing exactly
$n_d^{(m)}$ copies of each $d$ is the multinomial coefficient
\[
 \frac{m!}{\prod_{d\in D}n_d^{(m)}!}.
\]
By \eqref{eq:multinomial-estimate}, its logarithm is at least
\[
 -m\sum_{d\in D}q_d^{(m)}\log q_d^{(m)}
 -\log\binom{m+\card D-1}{\card D-1}
 =mG(p)+o(m).
\]
Indeed, $q_d^{(m)}\to p_d$ for every $d$, the function
$t\mapsto-t\log t$ is continuous on $[0,1]$ when its value at zero is
taken to be zero, and the final logarithmic term is $o(m)$ because $D$ is
fixed.  Also, for all sufficiently
large $m$,
\[
 \sum_{d\in D}n_d^{(m)}\kappa(d)\le(\mu+\varepsilon)m.
\]
More explicitly, each type count differs from $mp_d$ by $O_D(1)$, so the
rounding changes the total cost by $O_{D,\kappa}(1)$ and the normalized
entropy by $o(1)$.
Set
\[
 L_m=\left\lfloor\frac{(\mu+\varepsilon)m}{2}\right\rfloor
\]
and define
\begin{equation}\label{eq:Um}
 U_m=\left\{
   \sum_{i=0}^{m-1}a_iQ^i:
   a_i\in M,\ \sum_{i=0}^{m-1}a_i\le L_m
 \right\}.
\end{equation}

Take any sequence $(d_0,\ldots,d_{m-1})$ containing exactly
$n_d^{(m)}$ copies of each $d\in D$. Symmetry of the counts gives
$\sum_i d_i=0$. Hence
\[
 \sum_i\alpha_{d_i}-\sum_i\beta_{d_i}=0,
 \qquad
 \sum_i\alpha_{d_i}+\sum_i\beta_{d_i}=\sum_i\kappa(d_i),
\]
so each of the first two sums equals one half of the last. Each is an
integer not exceeding $(\mu+\varepsilon)m/2$, and hence is at most $L_m$.
Therefore
\[
 u=\sum_{i=0}^{m-1}\alpha_{d_i}Q^i,
 \qquad
 v=\sum_{i=0}^{m-1}\beta_{d_i}Q^i
\]
belong to $U_m$, and
\[
 u-v=\sum_{i=0}^{m-1}d_iQ^i.
\]
The map from difference-digit sequences to integers is injective. Indeed, if
\[
 \sum_i d_iQ^i=\sum_i d'_iQ^i,
\]
then, since $D\subseteq[-B,B]$, with $\delta_i=d_i-d'_i$ we have
$\sum_i\delta_iQ^i=0$ and $|\delta_i|\le2B<Q$. Reduction modulo
$Q=2B+1$ therefore gives $\delta_0=0$; after division by $Q$, induction
gives $\delta_i=0$ for every $i$. The multinomial estimate above therefore yields
\begin{equation}\label{eq:diff-lower}
 \liminf_{m\to\infty}\frac1m\log\card{U_m-U_m}\ge G(p).
\end{equation}
Every element of $U_m+U_m$ has a carry-free base-$Q$ digit sequence
$(s_0,\ldots,s_{m-1})\in S^m$ with $\sum_i s_i\le2L_m$. If $N_m$ is the
number of all such digit sequences $(s_0,\ldots,s_{m-1})\in S^m$ satisfying
$\sum_i s_i\le2L_m$, then, because $0<x<1$,
\[
 N_mx^{2L_m}
 \le\sum_{(s_i)\in S^m}x^{\sum_i s_i}
 =P_+(x)^m.
\]
Thus
\begin{equation}\label{eq:sum-upper}
 \limsup_{m\to\infty}\frac1m\log\card{U_m+U_m}
 \le \log P_+(x)-(\mu+\varepsilon)\log x.
\end{equation}
Combining \eqref{eq:diff-lower}, the identity defining $G(p)$, and
\eqref{eq:sum-upper} gives
\begin{equation}\label{eq:ratio-rate}
 \liminf_{m\to\infty}\frac1m
 \log\frac{\card{U_m-U_m}}{\card{U_m+U_m}}
 \ge R+\varepsilon\log x>0.
\end{equation}

Every $u\in U_m$ has digits at most $B$, so
\[
 u\le B\sum_{i=0}^{m-1}Q^i
   =\frac{Q^m-1}{2}.
\]
On the other hand, $\mu>0$ because $B\in D$ and $p_B>0$, so
$L_m\ge B$ for all sufficiently large $m$;
then $BQ^{m-1}\in U_m$.  Consequently
\[
 2(Q-1)Q^{m-1}+1
 \le 4\max U_m+1
 \le 2Q^m-1
\]
for all sufficiently large $m$, and hence
\[
 \lim_{m\to\infty}\frac1m\log(4\max U_m+1)=\log Q.
\]
The numerator in \eqref{eq:GHR-dilated} is positive for all sufficiently
large $m$ by \eqref{eq:ratio-rate}.  Apply \eqref{eq:GHR-dilated} to $U_m$
and take the lower limit of the resulting right-hand side.  The preceding
denominator limit and \eqref{eq:ratio-rate} give
\[
 C_{3a}\ge1+\frac{R+\varepsilon\log x}{\log Q}.
\]
For each fixed admissible $\varepsilon$, this conclusion follows after the
construction-size limit $m\to\infty$.  Since the left side is independent of
$\varepsilon$, it is at least the supremum of these lower bounds.  Letting
$\varepsilon\downarrow0$ in that supremum proves \eqref{eq:masked-bound}; no
interchange of two limits is being used.
\end{proof}

\begin{remark}[Recovery of Zheng's optimization]
For the complete mask $M=\{0,1,\ldots,B\}$,
\[
 P_+(x)=1+x+\cdots+x^{2B},
 \qquad
 P_-(x)=1+2x+2x^2+\cdots+2x^B.
\]
Maximizing Proposition~\ref{prop:masked} for $B=5$ reproduces Zheng's value
$1.173077279785\ldots$ \cite{Zheng2025}. Thus the proposition is a direct mask
generalization of Zheng's limiting construction, although the proof above is
self-contained apart from the finite-set principle \eqref{eq:GHR}
and the standard multinomial estimate \eqref{eq:multinomial-estimate}.
\end{remark}

\begin{remark}[The balancing condition at an interior optimum]\label{rem:moment}
Parameterize $x=e^{-\lambda}$, with $\lambda>0$, and write
\[
 F(\lambda)=\log P_-(e^{-\lambda})-\log P_+(e^{-\lambda}).
\]
Give $S$ probability weights proportional to $e^{-\lambda s}$ and give
$D$ probability weights proportional to $e^{-\lambda\kappa(d)}$.  Direct
differentiation gives
\[
 F'(\lambda)
 =\mathbb E_+[s]-\mathbb E_-[\kappa(d)].
\]
Consequently, every interior stationary point satisfies
\[
 \mathbb E_+[s]=\mathbb E_-[\kappa(d)].
\]
This identity is useful when comparing masks: the optimizing parameter
balances the typical weighted sum digit against the typical minimum
difference-witness cost.  It is not needed for the proof, which uses only
one fixed value of $\lambda$.
\end{remark}

\section{A carry-free masked-digit example}
Let
\begin{equation}\label{eq:record-mask}
 H=\angles{1518,1524,1587,2024,2032,2116}
\end{equation}
and let
\[
 M=H\cap[0,17032].
\]
Exact integer calculation gives
\[
 B=\max M=17032,
 \qquad Q=2B+1=34065,
\]
and
\begin{equation}\label{eq:counts}
 \card M=3121,\qquad
 \card{M+M}=18730,\qquad
 \card{M-M}=32369.
\end{equation}
The conductor of $H$ is $28922$.  Moreover,
\[
 [-13066,13066]\cap\mathbb Z\subseteq M-M,
\]
and the cheapest witness for difference one
has cost
\[
 \kappa(1)=14351.
\]
The verification package contains the complete mask, complete sum support,
and all values $(d,\kappa(d))$ for $d\in M-M$.

At the fixed decimal value
\[
 \lambda_0=
 0.000321149844434550835903084464712287774157626054669874186964,
 \qquad x_0=e^{-\lambda_0},
\]
the partition sums have decimal expansions
\[
 P_-(x_0)=668.3247325490738970446\ldots,
 \qquad
 P_+(x_0)=91.03099370518693032667\ldots,
\]
so
\[
 \log\frac{P_-(x_0)}{P_+(x_0)}
 =1.99357414328979926593\ldots
\]
and
\begin{equation}\label{eq:theta-numerical}
 1+\frac{\log(P_-(x_0)/P_+(x_0))}{\log 34065}
 =1.19102809750994590261\ldots.
\end{equation}
The accompanying MPFR verifier \cite{FousseEtAl2007} uses $384$-bit arithmetic with directed
rounding. It reconstructs all discrete data from the six generators and
proves the conservative strict inequality
\begin{equation}\label{eq:certified}
\boxed{C_{3a}>1.19102809}\qquad\text{(carry-free limit).}
\end{equation}
No optimization claim is needed for the proof: Proposition~\ref{prop:masked}
holds at the single fixed value $\lambda_0$.

\section{Controlled carries}
We now modify the masked-digit construction by replacing the carry-free base
$Q=2B+1$ by an arbitrary integer base $q>B$.  The elements under
consideration still have digits in $M\subseteq\{0,\ldots,B\}$ and a bound on
their total digit sum.  Their digitwise sums have digits in $M+M$, as before,
so the partition polynomial $P_+$ continues to give the required sumset
upper bound.  Signed difference digits lie in $D=M-M\subseteq[-B,B]$, but
when $q<2B+1$ distinct signed digit words can determine the same integer
after carrying.  We must count distinct differences without counting such
words more than once.

Fix a positive integer $k$ and group the digit positions into consecutive
blocks of $k$ positions.  Within one block, a signed digit vector
$(d_0,\ldots,d_{k-1})\in D^k$ has integer value
$\sum_{i=0}^{k-1}d_iq^i$.  Two such values have the same residue modulo
$q^k$ precisely when they differ by a carry into the next block.  We
therefore retain, for each residue class, the least possible witness cost of
a signed block with that residue.  This removes all within-block carry
collisions while preserving a large collection of differences.  Formally,
for $r\in\mathbb Z/q^k\mathbb Z$, define
\begin{equation}\label{eq:controlled-cost}
 c_k(r)=\min\left\{\sum_{i=0}^{k-1}\kappa(d_i):
 d_i\in D,\quad \sum_{i=0}^{k-1}d_iq^i\equiv r\pmod{q^k}\right\},
\end{equation}
with value $+\infty$ if no such digit vector exists.  Thus a vector
\emph{represents $r$ modulo $q^k$} when its integer value is congruent to
$r$ modulo $q^k$.  Put
\begin{equation}\label{eq:controlled-Z}
 Z_k(x)=\sum_{r\in\mathbb Z/q^k\mathbb Z}x^{c_k(r)},
 \qquad x^{+\infty}:=0.
\end{equation}

\begin{proposition}[Controlled-carry block bound]\label{prop:controlled-block}
For every $k\ge1$ and $x\in(0,1)$,
\begin{equation}\label{eq:controlled-block-bound}
 C_{3a}\ge1+
 \frac{k^{-1}\log Z_k(x)-\log P_+(x)}{\log q}.
\end{equation}
\end{proposition}

\begin{proof}
If the right side is at most $1$, use the elementary bound $C_{3a}\ge1$.
Otherwise fix $k$ for the remainder of this proof, let
\[
 \mathcal A_k=\{r\in\mathbb Z/q^k\mathbb Z:c_k(r)<+\infty\}
\]
be the attainable residue set, and set
\[
 R_k=\log Z_k(x)-k\log P_+(x)>0,
 \quad p_r=\frac{x^{c_k(r)}}{Z_k(x)}\quad(r\in\mathcal A_k),
 \quad \mu=\sum_{r\in\mathcal A_k}p_rc_k(r).
\]
Negation of all difference digits shows that $c_k(-r)=c_k(r)$ and hence
$p_{-r}=p_r$.  The entropy is
\begin{equation}\label{eq:controlled-entropy}
 H(p)=\log Z_k(x)-\mu\log x.
\end{equation}

Call a residue $r$ \emph{self-inverse} if $r=-r$ in
$\mathbb Z/q^k\mathbb Z$, equivalently if $2r=0$.  The zero residue is
always self-inverse; when $q^k$ is even, the only other self-inverse residue
is $h=q^k/2$.  Let $\mathcal R_k$ contain one member of each two-element
orbit $\{r,-r\}\subseteq\mathcal A_k$ under negation.  For every
$r\in\mathcal R_k$, choose a minimum-cost digit vector $d(r)$ representing
$r$ and set $d(-r)=-d(r)$.  For zero use the all-zero vector.  If $h$ lies
in $\mathcal A_k$, choose either minimum-cost vector $d(h)$; its negative
has the same cost and also represents $h$, because $-h=h$ modulo $q^k$.

Let $\ell$ be the number of blocks in the construction.  Define integer
type counts explicitly by
\begin{align*}
 N_r^{(\ell)}=N_{-r}^{(\ell)}&=\lfloor\ell p_r\rfloor
       &&(r\in\mathcal R_k),\\
 N_h^{(\ell)}&=2\left\lfloor\frac{\ell p_h}{2}\right\rfloor
       &&\text{if $h\in\mathcal A_k$},\\
 N_0^{(\ell)}&=\ell-2\sum_{r\in\mathcal R_k}
                    \lfloor\ell p_r\rfloor-N_h^{(\ell)},
\end{align*}
where the terms involving $h$ are omitted when $h\notin\mathcal A_k$.
These counts sum to $\ell$ and are nonnegative.  Indeed, every nonzero count
was rounded downward, so the defining identity for the probabilities gives
$N_0^{(\ell)}\ge\ell p_0>0$.  Moreover,
\[
 \frac{N_r^{(\ell)}}{\ell}\longrightarrow p_r
 \qquad(\ell\to\infty, r\in\mathcal A_k).
\]
Since $k$ is fixed, the alphabet $\mathcal A_k$ is finite.  Consequently,
the total rounding error in the cost is bounded by a constant depending on
$k$ (and on the fixed data $M,q,x$), and its contribution after division by
$\ell$ tends to zero.  The multinomial estimate
\eqref{eq:multinomial-estimate} gives
\begin{equation}\label{eq:controlled-mult}
 \log\frac{\ell!}{\prod_{r\in\mathcal A_k}N_r^{(\ell)}!}
 =\ell H(p)+o_k(\ell),
\end{equation}
and, for every fixed $\varepsilon>0$ and all sufficiently large $\ell$,
\[
 \sum_{r\in\mathcal A_k}N_r^{(\ell)}c_k(r)
 \le(\mu+\varepsilon)\ell.
\]
When $h\in\mathcal A_k$, order its occurrences in each residue word from
left to right, starting with index zero, and orient the even-indexed
occurrences as $d(h)$ and the odd-indexed occurrences as $-d(h)$.  Thus
consecutive occurrences are paired with opposite orientations.  The
definition of $N_h^{(\ell)}$ makes the two counts equal.  For
every $r\in\mathcal R_k$, the equal counts and the convention
$d(-r)=-d(r)$ cancel their signed digit vectors; the two orientations at
$h$ cancel in the same way.  Consequently, the sum of all $k\ell$ unweighted
difference digits is zero.

Use the symmetric minimum witnesses $\alpha_d,\beta_d$ fixed in the proof of
Proposition~\ref{prop:masked}.  The total $\alpha$-digit sum and total
$\beta$-digit sum are then equal, and each is half the total cost.  With
\[
 L_\ell=\left\lfloor\frac{(\mu+\varepsilon)\ell}{2}\right\rfloor,
\]
let $U_\ell$ be the set of base-$q$ integers with $k\ell$ digits in $M$ and
total digit sum at most $L_\ell$.  Every selected block-residue word supplies
a difference of two members of $U_\ell$.

For clarity, if a chosen signed block has integer value
\[
 R_j=\sum_{i=0}^{k-1}d_{jk+i}q^i,
\]
then
\[
 |R_j|\le B\frac{q^k-1}{q-1}\le q^k-1.
\]
Hence $R_j\equiv r_j\pmod{q^k}$ and $0\le r_j<q^k$ imply
$R_j\in\{r_j,r_j-q^k\}$.  We now verify injectivity of the complete
encoding.  Put $Q=q^k$.  Suppose two oriented residue words give the same
integer difference.  Reduction modulo $Q$ determines their first canonical
residue, so the first residues agree.  The corresponding signed blocks also
agree as integers: for an ordinary residue the chosen block is fixed, while
for $h$ the sign is determined by the parity of the number of earlier
occurrences of $h$.  After subtracting this common first block and dividing
by $Q$, the same argument applies to the two suffixes, with the same parity
offset because the removed residues agree.  Induction recovers the entire
residue word.  Thus distinct block-residue words give distinct integer
differences.

The prefix dependence of the alternating rule is essential to this decoding
argument.  A rule assigning $d(h)$ to the first half of all occurrences and
$-d(h)$ to the second half depends on the total number of occurrences in the
whole word and need not be injective.  The accompanying Lean development
includes an explicit collision for that alternative rule.
Thus
\[
 \liminf_{\ell\to\infty}\frac1\ell\log\card{U_\ell-U_\ell}\ge H(p).
\]
Every member of $U_\ell+U_\ell$ is produced by some word in $S^{k\ell}$ of
total at most $2L_\ell$.  Uniqueness is unnecessary because these words are
used only for an upper bound.  Therefore
\[
 \limsup_{\ell\to\infty}\frac1\ell\log\card{U_\ell+U_\ell}
 \le k\log P_+(x)-(\mu+\varepsilon)\log x.
\]
The logarithmic ratio consequently has lower rate
$R_k+\varepsilon\log x$, which is positive if
$0<\varepsilon<R_k/(-\log x)$.

Moreover, $\mu>0$: the residue $B$ is attainable and has positive cost and
positive $p$-weight.  Hence $L_\ell\ge B$ eventually and
$Bq^{k\ell-1}\in U_\ell$.  Together with
$\max U_\ell\le B(q^{k\ell}-1)/(q-1)$, this gives
\[
 \lim_{\ell\to\infty}\frac1\ell\log(4\max U_\ell+1)=k\log q.
\]
For each fixed admissible $\varepsilon$, apply
\eqref{eq:GHR-dilated} after taking $\ell$ sufficiently large, exactly as in
Proposition~\ref{prop:masked}.  Taking the supremum of the resulting bounds
over $\varepsilon$ and then letting $\varepsilon\downarrow0$ proves
\eqref{eq:controlled-block-bound}.
\end{proof}

\begin{lemma}[Infinite-block pressure]\label{lem:controlled-pressure}
For fixed $x\in(0,1)$, the limit
\[
 \mathcal P(x)=\lim_{k\to\infty}\frac1k\log Z_k(x)
 =\sup_{k\ge1}\frac1k\log Z_k(x)
\]
exists, and
\begin{equation}\label{eq:controlled-limit-bound}
 C_{3a}\ge1+\frac{\mathcal P(x)-\log P_+(x)}{\log q}.
\end{equation}
\end{lemma}

\begin{proof}
Concatenating minimum-cost representatives of residues modulo $q^k$ and
$q^\ell$ shows that
\[
 c_{k+\ell}(r+q^ks)\le c_k(r)+c_\ell(s).
\]
The map $(r,s)\mapsto r+q^ks$ is a bijection modulo $q^{k+\ell}$, so
$Z_{k+\ell}(x)\ge Z_k(x)Z_\ell(x)$.  Thus $a_k=\log Z_k(x)$, with
$a_0=0$, is
superadditive.  For completeness, if $k=tm+r$ with $0\le r<m$, repeated
superadditivity gives $a_k\ge t a_m+a_r$.  Since the finitely many $a_r$
are bounded below, first letting $k\to\infty$ and then taking the supremum
over $m$ proves
\[
 \lim_{k\to\infty}\frac{a_k}{k}
 =\sup_{m\ge1}\frac{a_m}{m};
\]
this is Fekete's lemma in the present setting.

The order of quantifiers is important.  Proposition~\ref{prop:controlled-block}
first fixes $k$ and then lets the number $\ell$ of blocks tend to infinity;
its rounding constants are allowed to depend on $k$.  The proposition
therefore proves a valid lower bound for each individual integer $k$.  We now
take the supremum of this already-established family of bounds.  No estimate
uniform in $k$, and no interchange of the limits $\ell\to\infty$ and
$k\to\infty$, is required.  Using the displayed identity for
$\mathcal P(x)$ gives \eqref{eq:controlled-limit-bound}.
\end{proof}

\subsection{Exact pressure for distinct sums}
The factor $P_+(x)$ used above counts sum-digit words, even when several
words produce the same integer after carrying.  The following construction
combines all such representations before counting them.

Let $\mathcal H_+=\{0,1\}$ be the possible carries.  For a canonical output
digit $r\in\{0,\ldots,q-1\}$, a transition $i\to j$ is possible when
\[
 s+i=r+qj
 \qquad\text{for some }s\in S=M+M.
\]
Its cost is $s$.  For a canonical word
$\mathbf r=(r_0,\ldots,r_{\ell-1})$ and $i,j\in\mathcal H_+$, let
$c^+_\ell(i,\mathbf r,j)$ be the least total cost of a path from $i$ to
$j$ emitting $\mathbf r$, with value $+\infty$ if no such path exists.  Set
\begin{equation}\label{eq:sum-block-matrix}
 A^+_\ell(x)_{ij}
 =\sum_{\mathbf r\in\{0,\ldots,q-1\}^\ell}
   x^{c^+_\ell(i,\mathbf r,j)}.
\end{equation}
Let $W_\ell(x)$ be the sum, over the distinct integer outputs of $\ell$
sum digits, of $x$ raised to the minimum cost of producing that output.
A canonical output word together with its terminal carry determines the
integer output uniquely, so the row sum of $A^+_\ell(x)$ starting with carry
zero is $W_\ell(x)$.  Starting with carry one translates every represented
integer by one and leaves all input digits and costs unchanged.  Hence both
row sums are $W_\ell(x)$, and therefore
\begin{equation}\label{eq:sum-row-root}
 A^+_\ell(x)\mathbf1=W_\ell(x)\mathbf1,
 \qquad \rho(A^+_\ell(x))=W_\ell(x).
\end{equation}

\begin{proposition}[Two-sided carry pressure]\label{prop:two-pressure}
Define
\[
 \mathcal P_-(x)=\mathcal P(x)
 =\sup_{k\ge1}\frac1k\log Z_k(x)
\]
and
\[
 \mathcal P_+(x)=\inf_{\ell\ge1}\frac1\ell\log W_\ell(x).
\]
Then the second infimum is a limit, and
\begin{equation}\label{eq:ultimate-bound}
 C_{3a}\ge
 1+\frac{\mathcal P_-(x)-\mathcal P_+(x)}{\log q}.
\end{equation}
More explicitly, every $k,\ell\ge1$ gives
\begin{equation}\label{eq:finite-two-block-bound}
 C_{3a}\ge1+
 \frac{k^{-1}\log Z_k(x)-\ell^{-1}\log W_\ell(x)}{\log q}.
\end{equation}
If a finite accessible difference-state set $F$ and a positive vector $v$
satisfy $T_{x,F}v\ge cv$, where $T_{x,F}$ is the principal submatrix of the
difference carry operator $T_x$ defined in
Section~\ref{sec:carry-state-operator}, and if
$W_\ell(x)^{1/\ell}\le p$, then
\begin{equation}\label{eq:finite-two-pressure-certificate}
 C_{3a}\ge1+\frac{\log(c/p)}{\log q}.
\end{equation}
Finally, each finite sum block has the rigorous error estimate
\begin{equation}\label{eq:sum-pressure-error}
 0\le \frac1\ell\log W_\ell(x)-\mathcal P_+(x)
 \le\frac{\log2}{\ell}.
\end{equation}
Thus $\mathcal P_+(x)$ is the exact asymptotic minimum-cost pressure of
distinct sum outputs.
\end{proposition}

\begin{proof}
Split a canonical output word of length $r+s$ at the block boundary.  For
fixed boundary carries, the weight of a concatenated output is the maximum,
over the two possible internal carries, of the two products of minimum-cost
block weights.  The corresponding entry of the ordinary matrix product
$A^+_r(x)A^+_s(x)$ is the sum of those at most two products.  Consequently,
after summing over outputs and terminal carries,
\begin{equation}\label{eq:sum-almost-multiplicative}
 W_{r+s}(x)\le W_r(x)W_s(x)\le2W_{r+s}(x).
\end{equation}
The left inequality makes $\log W_\ell(x)$ subadditive.  Fekete's lemma and
\eqref{eq:sum-row-root} give
\[
 \lim_{\ell\to\infty}\frac1\ell\log W_\ell(x)
 =\inf_{\ell\ge1}\frac1\ell\log W_\ell(x)
 =\mathcal P_+(x).
\]

Fix $k$ and $\ell$, and take the construction length $N$ to be a common
multiple of them.  The difference argument in
Proposition~\ref{prop:controlled-block} supplies the same entropy lower
bound per $k$-digit block.  If both input digit sums are at most $L$, every
element of the sumset has minimum sum-digit cost at most $2L$.  Since
$0<x<1$,
\[
 \card{U+U}x^{2L}\le W_N(x)
 \le W_\ell(x)^{N/\ell}.
\]
The digit-budget term cancels the identical term in the difference entropy
calculation.  Dividing by $N$ and then taking the construction-size limit
proves \eqref{eq:finite-two-block-bound}.  Taking the supremum over $k$ and
the infimum over $\ell$ proves \eqref{eq:ultimate-bound}.

Proposition~\ref{prop:controlled-finite-graph} gives
$\mathcal P_-(x)\ge\log c$, while
$\mathcal P_+(x)\le\ell^{-1}\log W_\ell(x)\le\log p$; substituting these
two inequalities in \eqref{eq:ultimate-bound} proves
\eqref{eq:finite-two-pressure-certificate}.  Finally, iterating the right
inequality in \eqref{eq:sum-almost-multiplicative} gives
\[
 W_{n\ell}(x)\ge 2^{-(n-1)}W_\ell(x)^n.
\]
Taking logarithms, dividing by $n\ell$, and letting $n\to\infty$ proves the
upper error bound in \eqref{eq:sum-pressure-error}; the lower bound follows
from the definition of the infimum.
\end{proof}

\subsection{The carry-state transfer operator}\label{sec:carry-state-operator}
Put $w(d)=\kappa(d)$ for $d\in D$ and $w(d)=+\infty$ otherwise, and define
for $0\le r\le q$
\[
 A_r^+=w(r),\qquad A_r^-=w(q-r).
\]
Thus $A_r^+$ is the cost of signed digit $r$, while $A_r^-$ is the cost of
signed digit $r-q$; the latter equality uses the evenness of $w$.
Here $A_q^+=+\infty$ and $A_q^-=w(0)=0$.  For a canonical prefix residue
$R\in\{0,\ldots,q^j-1\}$ for which at least one signed word exists, define
$C_j^{(0)}(R)$ and
$C_j^{(1)}(R)$ to be the minimum costs of signed $j$-digit words whose
integer values are, respectively, $R$ and $R-q^j$; an impossible value has
cost $+\infty$.  These are the only two possible integer values with residue
$R$, because
\[
 \left|\sum_{i<j}d_iq^i\right|
 \le B\frac{q^j-1}{q-1}\le q^j-1.
\]
Let $m_j(R)=\min(C_j^{(0)}(R),C_j^{(1)}(R))$ and normalize by setting
\[
 u_0=C_j^{(0)}(R)-m_j(R),\qquad
 u_1=C_j^{(1)}(R)-m_j(R).
\]
At least one of $u_0,u_1$ is zero.  The state is
\[
 \delta=u_1-u_0\in\mathbb Z\cup\{-\infty,+\infty\},
\]
with the conventions $\delta=+\infty$ when $u_1=+\infty$ and
$\delta=-\infty$ when $u_0=+\infty$.  Thus $\delta$ records both which of
the two representatives is cheaper and their exact cost difference.  For
the empty prefix, $C_0^{(0)}(0)=0$ and $C_0^{(1)}(0)=+\infty$, so the initial
state is $+\infty$.
Conversely, $\delta$ uniquely determines the normalized pair: it is
$(0,\delta)$ when $0\le\delta<+\infty$, $(-\delta,0)$ when
$-\infty<\delta<0$, $(0,+\infty)$ when $\delta=+\infty$, and
$(+\infty,0)$ when $\delta=-\infty$.  Hence the update below depends only
on the current state and the next canonical digit.

On appending the next canonical digit $r\in\{0,\ldots,q-1\}$, set
\begin{equation}\label{eq:carry-transition}
 v_0=\min(u_0+A_r^+,u_1+A_{r+1}^+),\qquad
 v_1=\min(u_0+A_r^-,u_1+A_{r+1}^-).
\end{equation}
If at least one value is finite, let
\[
 a(\delta,r)=\min(v_0,v_1),\qquad
 \Phi(\delta,r)=v_1-v_0
\]
with the same extended convention.  The definitions of $A_q^+$ and $A_q^-$
make \eqref{eq:carry-transition} valid also for $r=q-1$, including the
borrow across the new boundary.  Induction on the target digits shows that
the sum of the increments $a$ along a target word is exactly the minimum
cost $c_k(R)$ of its terminal residue $R\in\mathbb Z/q^k\mathbb Z$.

Let $T_x$ be the reachable-state kernel
\begin{equation}\label{eq:carry-operator}
 T_x(\delta,\delta')=
 \sum_{\substack{0\le r<q\\\Phi(\delta,r)=\delta'}}x^{a(\delta,r)}.
\end{equation}
For vectors $u$ and $v$ indexed by the reachable states, write
$u\cdot v=\sum_s u_s v_s$.  Let $e_{+\infty}$ denote
the coordinate vector supported at the initial state $+\infty$, and let
$\mathbf1$ denote the vector whose coordinates are all one.  Starting from
$\delta=+\infty$, we then have the exact scalar path identity
\begin{equation}\label{eq:carry-path-identity}
 Z_k(x)=e_{+\infty}\cdot\bigl(T_x^k\mathbf1\bigr).
\end{equation}

\begin{proposition}[Finite accessible-subgraph certificate]
\label{prop:controlled-finite-graph}
Let $\mathcal S$ be the set of states reachable from $+\infty$.  Suppose
$F\subseteq\mathcal S$ is finite and there is an integer $h\ge0$ such that
\[
 (T_x^h)(+\infty,s)>0\qquad(s\in F).
\]
Let $T_{x,F}=(T_x(s,t))_{s,t\in F}$ be the principal submatrix indexed by
$F$, and let $\rho(T_{x,F})$ denote its spectral radius, that is, the largest
modulus of one of its complex eigenvalues.  Then
\[
 \mathcal P(x)\ge\log\rho(T_{x,F}).
\]
Here and below, $\log0=-\infty$ when the retained matrix is nilpotent.
If $c>0$, $v>0$, and $T_{x,F}v\ge cv$ componentwise, then
\begin{equation}\label{eq:controlled-collatz}
 C_{3a}\ge1+\frac{\log c-\log P_+(x)}{\log q}.
\end{equation}
\end{proposition}

\begin{proof}
Define $b_s=(T_x^h)(+\infty,s)>0$ for $s\in F$.  In the path expansion of
\eqref{eq:carry-path-identity}, retain only paths that reach a state of $F$
after $h$ steps and remain in $F$ for the next $n$ steps.  Since every path
weight is nonnegative,
\[
 Z_{h+n}(x)\ge
 b\cdot\bigl(T_{x,F}^n\mathbf1\bigr).
\]
No irreducibility assumption is needed.  Put $A=T_{x,F}$ and
$b_*=\min_{s\in F}b_s>0$.  For the induced matrix $1$-norm,
\[
 b\cdot\bigl(A^n\mathbf1\bigr)
 \ge b_*\,\mathbf1\cdot\bigl(A^n\mathbf1\bigr)
 \ge b_*\lVert A^n\rVert_1.
\]
The Jordan normal form of the finite matrix $A$ gives
$\lim_{n\to\infty}\lVert A^n\rVert_1^{1/n}=\rho(A)$: each Jordan block with
eigenvalue $\lambda$ contributes a polynomial in $n$ times $|\lambda|^n$.
Hence the preceding path restriction and the existence of the pressure limit
imply
$\mathcal P(x)\ge\log\rho(A)$ (with the case $\rho(A)=0$ trivial).

Finally, if $v>0$ and $Av\ge cv$, then $A^nv\ge c^nv$ for every $n$.
Writing $v_* =\min_s v_s$ and $v^*=\max_s v_s$ gives, for the induced
infinity norm,
\[
 \lVert A^n\rVert_\infty
 \ge c^n\frac{v_*}{v^*}.
\]
The same Jordan-form norm limit yields $\rho(A)\ge c$.  This is
the positive-vector form of the Collatz--Wielandt bound, proved here for the
possibly reducible matrix $A$.  Equation \eqref{eq:controlled-collatz} now
follows from Lemma~\ref{lem:controlled-pressure}.
\end{proof}

\subsection{The controlled-carry certificate}\label{sec:controlled-certificate}
Take
\begin{equation}\label{eq:controlled-mask}
 B=26972,\qquad q=27022,
\qquad
 M=\angles{1971,2016,2100,2628,2688,2800}\cap[0,B],
\end{equation}
and
\begin{equation}\label{eq:controlled-x}
 x=\mathtt{0x1.ffde827adc0fep-1}
 =0.9997444891905898\ldots.
\end{equation}
Exact reconstruction gives $\card M=5869$ and $\card{M+M}=31274$.
For $j=1,\ldots,8$, let
\[
 F_j=\{s\in\mathcal S:(T_x^j)(+\infty,s)>0\}.
\]
The verifier obtains
\[
 \card{F_1}=7451,\quad \card{F_2}=49009,\quad
 \card{F_3}=109851,\quad \card{F_4}=161577,
\]
\[
 \card{F_5}=210577,\quad \card{F_6}=258233,\quad
 \card{F_7}=305267,\quad \card{F_8}=351951.
\]
It checks a positive dyadic vector componentwise for $T_{x,F_8}$ and proves
\[
 \rho(T_{x,F_8})>582.117820.
\]
The file \texttt{certificate/controlled/pf8\_vector.hex} contains the vector
$v$, one coordinate for each state of $F_8$.  Each line gives the integer
state label followed by the corresponding positive binary64 value in
hexadecimal notation.  The verifier reconstructs $F_8$ independently,
checks that the labels occur in the required order, and then verifies the
componentwise inequality
\[
 T_{x,F_8}v>582.117820v.
\]
All powers in the difference matrix are rounded downward.  A row contains at
most $q$ multiplications and $q$ additions; the verifier uses the deliberately
larger relative-error allowance $4q\,2^{-52}$ in every Collatz inequality.
Thus the matrix calculation is one-sided rather than a heuristic eigenvalue
estimate.

On the sum side, directed dynamic programming constructs $W_{8192}(x)$ and
proves
\[
 W_{8192}(x)^{1/8192}<79.428331.
\]
Every multiplication and addition in this calculation is rounded upward.
The values are rescaled by exact powers of two to avoid overflow and
underflow.  The final comparison is made against a downward-rounded power of
the rational number $79428331/10^6$.  By
\eqref{eq:sum-pressure-error}, replacing the exact infinite sum pressure by
this finite block changes the resulting exponent by at most
\[
 \frac{\log2}{8192\log27022}<8.292\times10^{-6}.
\]
This estimate is diagnostic only: the displayed finite-block upper bound is
already rigorous and is what is used in the theorem.

It remains to compare the rational pressure bounds with the reported decimal.
The accompanying Python script uses only exact rational arithmetic.  It
range-reduces each logarithm by powers of two and encloses
$\log y=2\operatorname{atanh}((y-1)/(y+1))$ between a positive rational
partial sum and an explicit geometric tail.  It certifies
\[
 1+\frac{\log(582.117820/79.428331)}{\log27022}>1.19519192.
\]
Equations \eqref{eq:finite-two-pressure-certificate} and
\eqref{eq:controlled-mask} therefore prove
\begin{equation}\label{eq:controlled-certified}
 \boxed{C_{3a}>1.19519192}.
\end{equation}
\section{A direct finite-depth certificate}\label{sec:finite-certificate}
For a bound that does not invoke either limiting theorem, return to the
carry-free mask \eqref{eq:record-mask}, with $B=17032$ and
$Q=34065$, and set
\[
 x_f=0.9996789017186568,\qquad N=10^{13}.
\]
Let
\[
 W=\sum_{d\in D}x_f^{\kappa(d)},\qquad
 n_d=\left\lfloor\frac{N x_f^{\kappa(d)}}{W}\right\rfloor\quad(d>0),
 \qquad n_{-d}=n_d,
\]
and let $n_0=N-2\sum_{d>0}n_d$.  Define
\[
 L=\frac12\sum_{d\in D}n_d\kappa(d)
 =51500691976683641
\]
and the explicit finite set
\[
 U=\left\{\sum_{i=0}^{N-1}a_iQ^i:
 a_i\in M,\ \sum_i a_i\le L\right\}.
\]
The finite set to which \eqref{eq:GHR} is applied is $V=2U$, as in the
dilation argument leading to \eqref{eq:GHR-dilated}.  The elementary bound
for $U$ itself is only $\card{U-U}\le2\max U+1$; since strict inequality has
not been established, applying \eqref{eq:GHR} directly to $U$ would leave a
gap.  Passing to $V$ supplies the required strictness at the negligible cost
of the additional $\log2$ described in Section~1.
For each $d$, use the symmetric minimum-cost witnesses
$(\alpha_d,\beta_d)$ fixed in the proof of
Proposition~\ref{prop:masked}.  There are exactly
\[
 D_N:=\frac{N!}{n_0!\prod_{d>0}(n_d!)^2}
\]
signed difference words $(d_0,\ldots,d_{N-1})$ having the prescribed counts.
Because $n_d=n_{-d}$, each such word satisfies $\sum_i d_i=0$.  Replacing
each $d_i$ by its fixed witness $\alpha_{d_i}-\beta_{d_i}$ gives two mask
words.  Their digit sums are equal because
$\sum_i(\alpha_{d_i}-\beta_{d_i})=\sum_i d_i=0$, and their sum is
$\sum_i(\alpha_{d_i}+\beta_{d_i})$.  Consequently, each digit sum is
\[
 \frac12\sum_i\bigl(\alpha_{d_i}+\beta_{d_i}\bigr)
 =\frac12\sum_{d\in D}n_d\kappa(d)=L.
\]
They therefore encode two members of $U$, whose difference is
$\sum_i d_iQ^i$.  Finally, $d_i\in[-B,B]$ and $Q=2B+1$, so the carry-free
injectivity established in Proposition~\ref{prop:masked} implies that two
different signed words produce two different integers $\sum_i d_iQ^i$.
Thus these $D_N$ words give $D_N$ distinct elements of $U-U$, and hence
\[
 \card{U-U}\ge D_N,
\]
while the same elementary weighted count used earlier gives
\[
 \card{U+U}\le S_N:=P_+(x_f)^N x_f^{-2L}.
\]
No limit is involved in either inequality.  Since
$\max U\le(Q^N-1)/2$, the dilated denominator satisfies
\[
 4\max U+1\le2Q^N-1<2Q^N.
\]
Directed $512$-bit MPFR evaluation \cite{FousseEtAl2007} reconstructs all
type counts (including $n_0=14962802398$) and certifies the strict inequalities
\begin{align*}
 \log D_N&>98126619915169.7944,\\
 \log S_N&<78190878820128.0915,\\
 \log(2Q^N)&<104360257432078.3785.
\end{align*}
Substitution in \eqref{eq:GHR-dilated} gives a lower value
$1.1910280942725\ldots$ and proves the wholly finite statement
\begin{equation}\label{eq:finite-certified}
 \boxed{C_{3a}>1.19102809}\qquad\text{(explicit finite depth $N=10^{13}$).}
\end{equation}
This value is weaker than both limiting bounds, as expected, but it can be
recomputed directly without accepting the asymptotic arguments in
Proposition~\ref{prop:masked} or Lemma~\ref{lem:controlled-pressure}.

\section{Semigroup structure and computational motivation}\label{sec:structure}
This section explains the numerical-semigroup structure of the carry-free
example in Section~3.  It does not assert that a mask optimal for
the carry-free base remains optimal after carries are allowed: the
controlled-carry certificate uses the different mask
\[
 \angles{1971,2016,2100,2628,2688,2800}\cap[0,26972]
 =\bigl(\angles{\{3,4\}\times\{657,672,700\}}\bigr)\cap[0,26972].
\]
Nevertheless, both successful masks arise from a $2\times3$ product grid,
which suggests that the relation structure discussed below remains useful
when the base is lowered even though the optimal parameters change.

A numerical semigroup is a cofinite additive submonoid of $\N$.  Numerical
semigroups are useful here because their missing elements suppress the
low-degree part of $P_+$, while additive relations among their generators
create further sum collisions.  The objective is nevertheless not determined
by the Hilbert series or by the three cardinalities in
\eqref{eq:counts}.  It also depends on the complete weighted minimum-cost
function $\kappa$.

For a generating vector $G=(g_1,\ldots,g_e)$, a useful untruncated analogue
of $\kappa$ is the residual cost function
\[
 \nu_G(d)=\min\left\{
 \sum_i(z_i^++z_i^-)g_i:
 z_i^\pm\in\N,\quad
 \sum_i(z_i^+-z_i^-)g_i=d
 \right\}.
\]
Thus $\nu_G(d)$ is the least total cost of two semigroup elements whose
difference is $d$.  If a minimizing pair lies in $[0,B]$, then
$\kappa(d)=\nu_G(d)$.  Exact relations, with residual zero, help produce
sum collisions; near relations with small nonzero residuals help produce
many difference digits at low cost.  The example combines both kinds.

\subsection{A progression of carry-free constructions}
The carry-free example emerged from a sequence of increasingly structured
numerical-semigroup masks:
\[
\begin{array}{c@{\quad}c@{\quad}c}
\toprule
\text{semigroup generators} & B & \text{resulting value}\\
\midrule
\{24,26,36,39\} & 189 & 1.1893936243\ldots\\
\{60,63,80,84\} & 538 & 1.1902074763\ldots\\
\{312,315,336,416,420\} & 3084 & 1.1902381380\ldots\\
\{1518,1524,1587,2024,2032,2116\} & 17032
 & 1.1910280975\ldots\\
\bottomrule
\end{array}
\]
The first two generator sets have the forms
\[
 12\angles{2,3}+13\angles{2,3}
 \quad\text{and}\quad
 20\angles{3,4}+21\angles{3,4}.
\]
In the general four-generator family
$S\angles{k,k+1}+(S+1)\angles{k,k+1}$, the identity
\[
 Sk+(S+1)(k+1)=S(k+1)+(S+1)k+1
\]
produces a unit residual with shared support.

The former five-generator example is an iterated gluing:
\[
 \angles{312,315,336,416,420}
 =21\angles{15,16,20}+104\angles{3,4},
\qquad
 \angles{15,16,20}=5\angles{3,4}+16\N.
\]
Gluing is standard in the theory of complete-intersection numerical
semigroups \cite{AssiGarciaSanchez2013,AssiGarciaOjeda2015}.
Studying this decomposition led to searches parameterized by relations
rather than by unrelated generator coordinates.

\subsection{A balanced simple gluing in the columns}
Write
\[
 C=\angles{506,508,529}.
\]
The column semigroup has the exact simple-gluing description
\begin{equation}\label{eq:column-gluing}
 C=23\angles{22,23}+508\N,
 \qquad 508=21\cdot22+2\cdot23.
\end{equation}
More generally, for a simple gluing
\[
 C_{a,b,r,s}=a\angles{r,s}+b\N
             =\angles{ar,as,b},
\]
where $\gcd(r,s)=\gcd(a,b)=1$ and $b\in\angles{r,s}$,
the two complete-intersection relation degrees are $ars$ and $ab$, and
its Hilbert series is
\[
 \sum_{n\in C_{a,b,r,s}}t^n
 =\frac{(1-t^{ars})(1-t^{ab})}
 {(1-t^{ar})(1-t^{as})(1-t^b)}.
\]
For these parameters
\[
 (r,s,a,b)=(22,23,23,508),
\]
the relation degrees are
\[
 ars=11638,\qquad ab=11684,
\]
differing by only $46$, while the generators
\[
 ar=506,\qquad b=508,\qquad as=529
\]
are also clustered.  The two balances have distinct meanings:
$b\approx rs$ aligns the relation degrees, whereas $b\approx ar$ clusters
two generators.  Satisfying both naturally forces $a\approx s$; in the
present example, $a=s$ exactly.  Formula \eqref{eq:column-gluing} also gives
the column Frobenius number
\[
 23(22\cdot23-22-23)+22\cdot508=21779.
\]

\subsection{\texorpdfstring{The $2\times3$ generator grid}
                           {The two-by-three generator grid}}
The semigroup is generated by the grid
\begin{equation}\label{eq:grid}
 \{3,4\}\times\{506,508,529\}
 =
 \{1518,1524,1587,2024,2032,2116\}.
\end{equation}
For a column $c$, the two row generators have the exact relation
\[
 4(3c)=3(4c).
\]
For the three columns, these relations occur at degrees
\[
 6072,\qquad6096,\qquad6348,
\]
well before the cutoff $B=17032$.

The column relation $23\cdot506=22\cdot529$ also lifts to the grid:
\[
 5(1518)+2(2024)
 =2(1587)+4(2116)
 =11638.
\]
The other relation in \eqref{eq:column-gluing} can be lifted after adding
$4\cdot529$ to both sides:
\[
 5(1524)+2(2032)+2116
 =7(1518)+2(1587)
 =13800.
\]
Thus several exact relation layers are already active below the chosen
cutoff, while the conductor of the full six-generator semigroup,
$28922$, remains substantially larger.

The same grid simultaneously gives cheap near relations.  If
$c_i<c_j$ are columns, then
\begin{equation}\label{eq:grid-residual}
 (3c_i+4c_j)-(4c_i+3c_j)=c_j-c_i,
\end{equation}
and the total cost of the two witnesses is $7(c_i+c_j)$.  The column gaps
are $2$, $21$, and $23$, so the mask has the explicit bounds
\[
 \kappa(2)\le7098,\qquad
 \kappa(21)\le7259,\qquad
 \kappa(23)\le7245.
\]
Because $\gcd(2,21)=1$, the cheap residual directions generate the full
integer lattice, although combining them can increase witness cost and is
subject to the cutoff.  In contrast, the exact certificate gives
$\kappa(1)=14351$.  This illustrates why minimizing $\kappa(1)$ alone is
not an adequate design criterion: Proposition~\ref{prop:masked} uses the
entire weighted distribution of $\kappa(d)$.

\subsection{Structural properties suggested by the search}
Searches around \eqref{eq:grid} found a narrow interior ridge: consecutive
$r,s$, equality $a=s$, near equality $b\approx rs\approx ar$, and the row
pair $\{3,4\}$ were all favored.  Relaxing the relation balance and adding
a third binary product factor did not improve the best value found.  These
computations are evidence about the explored families, not an optimality
proof.

The structural evidence suggests that a strong semigroup mask should have:
\begin{enumerate}
 \item several exact relations becoming active below the cutoff;
 \item several cheap near relations with small coprime residuals;
 \item substantial sharing of generators among those relations;
 \item clustered relation degrees, so that multiple collision mechanisms
       become active in the same range; and
 \item delayed density, with the cutoff after the useful relation layers
       but before the conductor.
\end{enumerate}
Product grids and gluings manufacture these features, but neither is
logically required by Proposition~\ref{prop:masked}.  A natural broader
program is therefore to prescribe a small weighted relation lattice first
and solve the resulting Diophantine constraints for the generators.  This
includes low-rank grids with small circuit defects, carefully selected
additive gluings, and constrained families described by numerical-semigroup
coordinates.  Such a
relation-template approach is more closely tied to the objective than an
unstructured search over generator tuples.


\section{Verification and reproducibility}
The source code, certificate files, verification scripts, and saved run
logs are available in the repository at
\[
 \packageurl.
\]
After this update package is integrated, the repository contains the
following checks.  Items 1--3 are the existing carry-free checks; items
4--6 are supplied or extended by this update, and item 7 is the accompanying
formal verification.
\begin{enumerate}
 \item \texttt{verification/verify\_mpfr.cpp} reconstructs $M$, $M+M$, and
 $\kappa(d)$ using exact integer arithmetic and certifies
 \eqref{eq:certified} using MPFR directed rounding.
 \item \texttt{verification/verify\_python.py} independently reconstructs
 the exact discrete data using the Python standard library and checks the
 shipped certificate files and hashes.
 \item \texttt{verification/verify\_gp.gp} independently reconstructs the
 data in PARI/GP and evaluates the bound at high precision. This is a
 secondary numerical cross-check; the MPFR program supplies the rigorous
 floating-point certificate.
 \item \texttt{verification/controlled/certify\_candidate.cpp}, together
 with \texttt{pressure.hpp}, reconstructs the controlled-carry mask, all
 eight accessible state sets, the $351951$ component Collatz certificate,
 and the directed $8192$-digit sum-pressure bound.
 \item \texttt{verification/controlled/verify\_exponent.py} checks the final
 decimal exponent with exact rational logarithm enclosures.  Its conclusion
 does not depend on a floating-point implementation of $\log$.
 \item \texttt{verification/verify\_finite\_mpfr.cpp} reconstructs the
 $N=10^{13}$ finite type counts and certifies
 \eqref{eq:finite-certified} with directed MPFR rounding.  Its running time
 is governed by the fixed mask size, not by enumeration of the elements of
 $U$.
 \item \texttt{formalization/lean/} contains the Lean~4/Mathlib
 formalization produced with Aristotle.  Running its \texttt{verify.sh}
 builds the project, prints the principal theorem and axiom reports, and
 rejects any occurrence of \texttt{sorry}, \texttt{admit}, or an explicit
 project axiom.
\end{enumerate}

\begin{thebibliography}{99}
\bibitem{Csiszar1998}
I. Csisz\'ar,
\emph{The method of types},
IEEE Transactions on Information Theory 44 (1998), no.~6, 2505--2523,
\url{https://doi.org/10.1109/18.720546},
\url{https://iest2.ie.cuhk.edu.hk/~whyeung/post/manuscript/06_SAEP/method\%20of\%20types.pdf}.

\bibitem{OptimizationConstants2026}
D. Davis, P. Ivanisvili, T. Tao, and contributors,
\emph{Optimization Constants in Mathematics},
GitHub repository, 2026,
\url{https://github.com/teorth/optimizationproblems}.

\bibitem{GHR2007}
K. Gyarmati, F. Hennecart, and I. Z. Ruzsa,
\emph{Sums and differences of finite sets},
Functiones et Approximatio Commentarii Mathematici 37 (2007), 175--186,
\url{https://doi.org/10.7169/facm/1229618749},
\url{https://gyarmatikati.web.elte.hu/publ/sumdiffv.pdf}.

\bibitem{GeorgievEtAl2025}
B. Georgiev, J. G\'omez-Serrano, T. Tao, and A. Z. Wagner,
\emph{Mathematical exploration and discovery at scale},
arXiv:2511.02864, 2025,
\url{https://arxiv.org/abs/2511.02864}.

\bibitem{Gerbicz2025}
R. Gerbicz,
\emph{Sums and differences of sets (improvement over AlphaEvolve)},
arXiv:2505.16105, 2025,
\url{https://arxiv.org/abs/2505.16105}.

\bibitem{Zheng2025}
F. Zheng,
\emph{Sums and differences of sets: a further improvement over AlphaEvolve},
arXiv:2506.01896, 2025,
\url{https://arxiv.org/abs/2506.01896}.

\bibitem{Griego2026}
S. Griego,
\emph{Base-$21$ digit construction certificate for $C_{3a}$},
GitHub pull request \#71 to the Optimization Constants repository, 2026,
\url{https://github.com/teorth/optimizationproblems/pull/71}.

\bibitem{Mosaic2026}
Mosaic Intelligence,
\emph{Exact-count certificate for problem 3a},
GitHub pull request \#95 to the Optimization Constants repository, 2026,
\url{https://github.com/teorth/optimizationproblems/pull/95}.

\bibitem{Numaro2026}
Numaro,
\emph{Certified lower bound $C_{3a}\ge1.187326127925948$},
GitHub pull request \#131 to the Optimization Constants repository, 2026,
\url{https://github.com/teorth/optimizationproblems/pull/131}.

\bibitem{Kleinwaks2026}
L. Kleinwaks,
\emph{Masked-semigroup digit bound $C_{3a}>1.19102809$},
GitHub pull request \#134 to the Optimization Constants repository, 2026,
\url{https://github.com/teorth/optimizationproblems/pull/134}.

\bibitem{AssiGarciaSanchez2013}
A. Assi and P. A. Garc\'ia-S\'anchez,
\emph{Constructing the set of complete intersection numerical semigroups
with a given Frobenius number},
arXiv:1204.4258, 2013,
\url{https://arxiv.org/abs/1204.4258}.

\bibitem{AssiGarciaOjeda2015}
A. Assi, P. A. Garc\'ia-S\'anchez, and I. Ojeda,
\emph{Frobenius vectors, Hilbert series and gluings of affine semigroups},
Journal of Commutative Algebra 7 (2015), no.~3, 317--335,
\url{https://doi.org/10.1216/JCA-2015-7-3-317},
\url{https://arxiv.org/abs/1311.1988}.

\bibitem{FousseEtAl2007}
L. Fousse, G. Hanrot, V. Lef\`evre, P. P\'elissier, and P. Zimmermann,
\emph{MPFR: A multiple-precision binary floating-point library with correct
rounding},
ACM Transactions on Mathematical Software 33 (2007), no.~2, Article~13,
\url{https://doi.org/10.1145/1236463.1236468},
\url{https://inria.hal.science/inria-00070266v1/document}.

\bibitem{Aristotle2025}
T. Achim et al.,
\emph{Aristotle: IMO-level Automated Theorem Proving},
arXiv:2510.01346, 2025,
\url{https://arxiv.org/abs/2510.01346}.

\end{thebibliography}

\end{document}
